\chapter{Introducción}

La programación competitiva combina dos temas:
(1) el diseño de algoritmos y (2) la implementación de algoritmos.

El \key{diseño de algoritmos} consiste en la resolución de problemas
y el pensamiento matemático.
Se requieren habilidades para analizar problemas y resolverlos
de forma creativa.
Un algoritmo para resolver un problema
tiene que ser a la vez correcto y eficiente,
y el núcleo del problema suele consistir
en idear un algoritmo eficiente.

El conocimiento teórico de los algoritmos
es importante para los programadores competitivos.
Por lo general, la solución de un problema es
una combinación de técnicas bien conocidas e
ideas nuevas.
Las técnicas que aparecen en la programación competitiva
también constituyen la base de la investigación
científica de los algoritmos.

La \key{implementación de algoritmos} requiere buenas
habilidades de programación.
En la programación competitiva, las soluciones
se califican probando un algoritmo implementado
mediante un conjunto de casos de prueba.
Por lo tanto, no basta con que la idea del
algoritmo sea correcta, sino que la implementación también
tiene que ser correcta.

Un buen estilo de código en los concursos es
directo y conciso.
Los programas deben escribirse con rapidez,
porque no hay mucho tiempo disponible.
A diferencia de la ingeniería de software tradicional,
los programas son cortos (normalmente a lo sumo unos pocos
cientos de líneas de código), y no necesitan
mantenerse después del concurso.

\section{Lenguajes de programación}

\index{lenguaje de programación}

En la actualidad, los lenguajes de programación
más populares usados en los concursos son C++, Python y Java.
Por ejemplo, en el Google Code Jam 2017,
entre los mejores 3.000 participantes,
el 79 \% usó C++,
el 16 \% usó Python y
el 8 \% usó Java \cite{goo17}.
Algunos participantes también usaron varios lenguajes.

Mucha gente piensa que C++ es la mejor opción
para un programador competitivo,
y C++ casi siempre está disponible en
los sistemas de concursos.
Los beneficios de usar C++ son que
es un lenguaje muy eficiente y
su biblioteca estándar contiene una
gran colección
de estructuras de datos y algoritmos.

Por otro lado, es bueno
dominar varios lenguajes y comprender
sus fortalezas.
Por ejemplo, si se necesitan enteros muy grandes
en el problema,
Python puede ser una buena opción, porque
contiene operaciones integradas para
realizar calculos con enteros grandes.
Aun así, la mayoría de los problemas en los concursos de programación
se plantean de modo que usar un lenguaje de programación específico
no sea una ventaja injusta.

Todos los programas de ejemplo de este libro están escritos en C++,
y a menudo se usan las estructuras de datos y los algoritmos de la biblioteca estándar.
Los programas siguen el estándar C++11,
que puede usarse en la mayoría de los concursos hoy en día.
Si todavía no sabes programar en C++,
ahora es un buen momento para empezar a aprender.

\subsubsection{Plantilla de código C++}

Una plantilla de código C++ típica para programación competitiva
tiene este aspecto:

\begin{lstlisting}
#include <bits/stdc++.h>

using namespace std;

int main() {
    // aquí se escribe la solución
}
\end{lstlisting}

La línea \texttt{\#include} al principio
del código es una instrucción que le indica al compilador \texttt{g++}
 incluir toda la biblioteca estándar.
Así, no es necesario incluir por separado
bibliotecas como \texttt{iostream},
\texttt{vector} y \texttt{algorithm},
sino que están disponibles automáticamente.

La línea \texttt{using} declara
que las clases y funciones
de la biblioteca estándar pueden usarse directamente
en el código.
Sin la línea \texttt{using} tendríamos
que escribir, por ejemplo, \texttt{std::cout},
pero ahora basta con escribir \texttt{cout}.

El código puede compilarse usando el siguiente comando:

\begin{lstlisting}
g++ -std=c++11 -O2 -Wall test.cpp -o test
\end{lstlisting}

Este comando produce un archivo binario \texttt{test}
a partir del código fuente \texttt{test.cpp}.
El compilador sigue el estándar C++11
(\texttt{-std=c++11}),
optimiza el código (\texttt{-O2})
y muestra advertencias sobre posibles errores (\texttt{-Wall}).

\section{Entrada y salida}

\index{entrada y salida}

En la mayoría de los concursos, se usan flujos estándar para
leer la entrada y escribir la salida.
En C++, los flujos estándar son
\texttt{cin} para la entrada y \texttt{cout} para la salida.
Además, pueden usarse las funciones de C
\texttt{scanf} y \texttt{printf}.

La entrada del programa normalmente consta de
números y cadenas separados por
espacios y saltos de línea.
Pueden leerse del flujo \texttt{cin}
de la siguiente manera:

\begin{lstlisting}
int a, b;
string x;
cin >> a >> b >> x;
\end{lstlisting}

Este tipo de código siempre funciona,
suponiendo que haya al menos un espacio
o salto de línea entre cada elemento de la entrada.
Por ejemplo, el código anterior puede leer
ambas de las siguientes entradas:
\begin{lstlisting}
123 456 monkey
\end{lstlisting}
\begin{lstlisting}
123    456
monkey
\end{lstlisting}
El flujo \texttt{cout} se usa para la salida
de la siguiente manera:
\begin{lstlisting}
int a = 123, b = 456;
string x = "monkey";
cout << a << " " << b << " " << x << "\n";
\end{lstlisting}

La entrada y la salida a veces son
un cuello de botella en el programa.
Las siguientes líneas al principio del código
hacen que la entrada y la salida sean más eficientes:

\begin{lstlisting}
ios::sync_with_stdio(0);
cin.tie(0);
\end{lstlisting}

Nota que el salto de línea \texttt{"\textbackslash n"}
funciona más rápido que \texttt{endl},
porque \texttt{endl} siempre provoca
una operación de vaciado (flush).

Las funciones de C \texttt{scanf}
y \texttt{printf} son una alternativa
a los flujos estándar de C++.
Suelen ser un poco más rápidas,
pero también son más difíciles de usar.
El siguiente código lee dos enteros de la entrada:
\begin{lstlisting}
int a, b;
scanf("%d %d", &a, &b);
\end{lstlisting}
El siguiente código imprime dos enteros:
\begin{lstlisting}
int a = 123, b = 456;
printf("%d %d\n", a, b);
\end{lstlisting}

A veces el programa debe leer una línea completa
de la entrada, que posiblemente contenga espacios.
Esto puede lograrse usando la
función \texttt{getline}:

\begin{lstlisting}
string s;
getline(cin, s);
\end{lstlisting}

Si la cantidad de datos es desconocida, el siguiente
bucle es útil:
\begin{lstlisting}
while (cin >> x) {
    // código
}
\end{lstlisting}
Este bucle lee elementos de la entrada
uno tras otro, hasta que no hay
más datos disponibles en la entrada.

En algunos sistemas de concursos, se usan archivos para
la entrada y la salida.
Una solución sencilla para esto es escribir
el código como de costumbre usando flujos estándar,
pero añadir las siguientes líneas al principio del código:
\begin{lstlisting}
freopen("input.txt", "r", stdin);
freopen("output.txt", "w", stdout);
\end{lstlisting}
Después de esto, el programa lee la entrada del archivo
''input.txt'' y escribe la salida en el archivo
''output.txt''.

\section{Trabajar con números}

\index{entero}

\subsubsection{Enteros}

El tipo entero más usado en programación competitiva
es \texttt{int}, que es un tipo de 32 bits con
un rango de valores de $-2^{31} \ldots 2^{31}-1$
o aproximadamente $-2 \cdot 10^9 \ldots 2 \cdot 10^9$.
Si el tipo \texttt{int} no es suficiente,
puede usarse el tipo de 64 bits \texttt{long long}.
Tiene un rango de valores de $-2^{63} \ldots 2^{63}-1$
o aproximadamente $-9 \cdot 10^{18} \ldots 9 \cdot 10^{18}$.

El siguiente código define una
variable \texttt{long long}:
\begin{lstlisting}
long long x = 123456789123456789LL;
\end{lstlisting}
El sufijo \texttt{LL} significa que el
tipo del número es \texttt{long long}.

Un error común al usar el tipo \texttt{long long}
es que el tipo \texttt{int} se sigue usando en algún lugar
del código.
Por ejemplo, el siguiente código contiene
un error sutil:

\begin{lstlisting}
int a = 123456789;
long long b = a*a;
cout << b << "\n"; // -1757895751
\end{lstlisting}

Aunque la variable \texttt{b} es de tipo \texttt{long long},
ambos números en la expresión \texttt{a*a}
son de tipo \texttt{int} y el resultado
también es de tipo \texttt{int}.
Por eso, la variable \texttt{b}
contendrá un resultado erróneo.
El problema puede resolverse cambiando el tipo
de \texttt{a} a \texttt{long long} o
cambiando la expresión a \texttt{(long long)a*a}.

Normalmente los problemas de los concursos se plantean de modo que el
tipo \texttt{long long} sea suficiente.
Aun así, es bueno saber que
el compilador \texttt{g++} también proporciona
un tipo de 128 bits \texttt{\_\_int128\_t}
con un rango de valores de
$-2^{127} \ldots 2^{127}-1$ o aproximadamente $-10^{38} \ldots 10^{38}$.
Sin embargo, este tipo no está disponible en todos los sistemas de concursos.

\subsubsection{Aritmética modular}

\index{residuo}
\index{aritmética modular}

Denotamos por $x \bmod m$ el residuo
cuando $x$ se divide entre $m$.
Por ejemplo, $17 \bmod 5 = 2$,
porque $17 = 3 \cdot 5 + 2$.

A veces, la respuesta a un problema es un
número muy grande pero basta con
darla ''módulo $m$'', es decir,
el residuo cuando la respuesta se divide entre $m$
(por ejemplo, ''módulo $10^9+7$'').
La idea es que aunque la respuesta real
sea muy grande,
basta con usar los tipos
\texttt{int} y \texttt{long long}.

Una propiedad importante del residuo es que
en la suma, la resta y la multiplicación,
el residuo puede tomarse antes de la operación:

\[
\begin{array}{rcr}
(a+b) \bmod m & = & (a \bmod m + b \bmod m) \bmod m \\
(a-b) \bmod m & = & (a \bmod m - b \bmod m) \bmod m \\
(a \cdot b) \bmod m & = & (a \bmod m \cdot b \bmod m) \bmod m
\end{array}
\]

Así, podemos tomar el residuo después de cada operación
y los números nunca se harán demasiado grandes.

Por ejemplo, el siguiente código calcula $n!$,
el factorial de $n$, módulo $m$:
\begin{lstlisting}
long long x = 1;
for (int i = 2; i <= n; i++) {
    x = (x*i)%m;
}
cout << x%m << "\n";
\end{lstlisting}

Normalmente queremos que el residuo esté siempre
entre $0\ldots m-1$.
Sin embargo, en C++ y otros lenguajes,
el residuo de un número negativo
es cero o negativo.
Una forma sencilla de asegurarse de que no haya
residuos negativos es calcular primero
el residuo como de costumbre y luego sumar $m$
si el resultado es negativo:
\begin{lstlisting}
x = x%m;
if (x < 0) x += m;
\end{lstlisting}
Sin embargo, esto solo es necesario cuando hay
restas en el código y el
residuo puede volverse negativo.

\subsubsection{Números de punto flotante}

\index{número de punto flotante}

Los tipos de punto flotante habituales en
programación competitiva son
el \texttt{double} de 64 bits
y, como extensión en el compilador \texttt{g++},
el \texttt{long double} de 80 bits.
En la mayoría de los casos, \texttt{double} es suficiente,
pero \texttt{long double} es más preciso.

La precisión requerida de la respuesta
suele indicarse en el enunciado del problema.
Una forma sencilla de mostrar la respuesta es usar
la función \texttt{printf}
e indicar el número de posiciones decimales
en la cadena de formato.
Por ejemplo, el siguiente código imprime
el valor de $x$ con 9 posiciones decimales:

\begin{lstlisting}
printf("%.9f\n", x);
\end{lstlisting}

Una dificultad al usar números de punto flotante
es que algunos números no pueden representarse
con exactitud como números de punto flotante,
y habrá errores de redondeo.
Por ejemplo, el resultado del siguiente código
es sorprendente:

\begin{lstlisting}
double x = 0.3*3+0.1;
printf("%.20f\n", x); // 0.99999999999999988898
\end{lstlisting}

Debido a un error de redondeo,
el valor de \texttt{x} es un poco menor que 1,
mientras que el valor correcto sería 1.

Es arriesgado comparar números de punto flotante
con el operador \texttt{==},
porque es posible que los valores debieran ser
iguales pero no lo sean debido a errores de precisión.
Una mejor forma de comparar números de punto flotante
es suponer que dos números son iguales
si la diferencia entre ellos es menor que $\varepsilon$,
donde $\varepsilon$ es un número pequeño.

En la práctica, los números pueden compararse
de la siguiente manera ($\varepsilon=10^{-9}$):

\begin{lstlisting}
if (abs(a-b) < 1e-9) {
    // a and b are equal
}
\end{lstlisting}

Nota que aunque los números de punto flotante son inexactos,
los enteros hasta cierto límite todavía pueden
representarse con exactitud.
Por ejemplo, usando \texttt{double},
es posible representar con exactitud todos los
enteros cuyo valor absoluto sea a lo sumo $2^{53}$.

\section{Acortar el código}

El código corto es ideal en programación competitiva,
porque los programas deben escribirse
lo más rápido posible.
Por eso, los programadores competitivos a menudo definen
nombres más cortos para los tipos de datos y otras partes del código.

\subsubsection{Nombres de tipos}
\index{tuppdef@\texttt{typedef}}
Usando el comando \texttt{typedef}
es posible dar un nombre más corto
a un tipo de dato.
Por ejemplo, el nombre \texttt{long long} es largo,
así que podemos definir un nombre más corto \texttt{ll}:
\begin{lstlisting}
typedef long long ll;
\end{lstlisting}
Después de esto, el código
\begin{lstlisting}
long long a = 123456789;
long long b = 987654321;
cout << a*b << "\n";
\end{lstlisting}
puede acortarse de la siguiente manera:
\begin{lstlisting}
ll a = 123456789;
ll b = 987654321;
cout << a*b << "\n";
\end{lstlisting}

El comando \texttt{typedef}
también puede usarse con tipos más complejos.
Por ejemplo, el siguiente código da
el nombre \texttt{vi} a un vector de enteros
y el nombre \texttt{pi} a un par
que contiene dos enteros.
\begin{lstlisting}
typedef vector<int> vi;
typedef pair<int,int> pi;
\end{lstlisting}

\subsubsection{Macros}
\index{macro}
Otra forma de acortar el código es definir
\key{macros}.
Una macro significa que ciertas cadenas del
código se cambiarán antes de la compilación.
En C++, las macros se definen usando la
palabra clave \texttt{\#define}.

Por ejemplo, podemos definir las siguientes macros:
\begin{lstlisting}
#define F first
#define S second
#define PB push_back
#define MP make_pair
\end{lstlisting}
Después de esto, el código
\begin{lstlisting}
v.push_back(make_pair(y1,x1));
v.push_back(make_pair(y2,x2));
int d = v[i].first+v[i].second;
\end{lstlisting}
puede acortarse de la siguiente manera:
\begin{lstlisting}
v.PB(MP(y1,x1));
v.PB(MP(y2,x2));
int d = v[i].F+v[i].S;
\end{lstlisting}

Una macro también puede tener parámetros
lo que hace posible acortar bucles y otras
estructuras.
Por ejemplo, podemos definir la siguiente macro:
\begin{lstlisting}
#define REP(i,a,b) for (int i = a; i <= b; i++)
\end{lstlisting}
Después de esto, el código
\begin{lstlisting}
for (int i = 1; i <= n; i++) {
    search(i);
}
\end{lstlisting}
puede acortarse de la siguiente manera:
\begin{lstlisting}
REP(i,1,n) {
    search(i);
}
\end{lstlisting}

A veces las macros provocan errores que pueden ser difíciles
de detectar. Por ejemplo, considera la siguiente macro
que calcula el cuadrado de un número:
\begin{lstlisting}
#define SQ(a) a*a
\end{lstlisting}
Esta macro \emph{no} siempre funciona como se espera.
Por ejemplo, el código
\begin{lstlisting}
cout << SQ(3+3) << "\n";
\end{lstlisting}
corresponde al código
\begin{lstlisting}
cout << 3+3*3+3 << "\n"; // 15
\end{lstlisting}

Una mejor versión de la macro es la siguiente:
\begin{lstlisting}
#define SQ(a) (a)*(a)
\end{lstlisting}
Ahora el código
\begin{lstlisting}
cout << SQ(3+3) << "\n";
\end{lstlisting}
corresponde al código
\begin{lstlisting}
cout << (3+3)*(3+3) << "\n"; // 36
\end{lstlisting}


\section{Matemáticas}

Las matemáticas desempeñan un papel importante en la programación
competitiva, y no es posible llegar a ser
un programador competitivo exitoso sin
tener buenas habilidades matemáticas.
Esta sección analiza algunos conceptos y fórmulas
matemáticos importantes que
se necesitarán más adelante en el libro.

\subsubsection{Fórmulas de sumas}

Cada suma de la forma
\[\sum_{x=1}^n x^k = 1^k+2^k+3^k+\ldots+n^k,\]
donde $k$ es un entero positivo,
tiene una fórmula cerrada que es un
polinomio de grado $k+1$.
Por ejemplo\footnote{\index{fórmula de Faulhaber}
Existe incluso una fórmula general para tales sumas, llamada \key{fórmula de Faulhaber},
pero es demasiado compleja para presentarla aquí.},
\[\sum_{x=1}^n x = 1+2+3+\ldots+n = \frac{n(n+1)}{2}\]
y
\[\sum_{x=1}^n x^2 = 1^2+2^2+3^2+\ldots+n^2 = \frac{n(n+1)(2n+1)}{6}.\]

Una \key{progresión aritmética} es una \index{progresión aritmética}
sucesión de números
donde la diferencia entre dos números consecutivos
cualesquiera es constante.
Por ejemplo,
\[3, 7, 11, 15\]
es una progresión aritmética con constante 4.
La suma de una progresión aritmética puede calcularse
usando la fórmula
\[\underbrace{a + \cdots + b}_{n \,\, \textrm{números}} = \frac{n(a+b)}{2}\]
donde $a$ es el primer número,
$b$ es el último número y
$n$ es la cantidad de números.
Por ejemplo,
\[3+7+11+15=\frac{4 \cdot (3+15)}{2} = 36.\]
La fórmula se basa en el hecho
de que la suma consta de $n$ números y
el valor de cada número es $(a+b)/2$ en promedio.

\index{progresión geométrica}
Una \key{progresión geométrica} es una sucesión
de números
donde la razón entre dos números consecutivos
cualesquiera es constante.
Por ejemplo,
\[3,6,12,24\]
es una progresión geométrica con constante 2.
La suma de una progresión geométrica puede calcularse
usando la fórmula
\[a + ak + ak^2 + \cdots + b = \frac{bk-a}{k-1}\]
donde $a$ es el primer número,
$b$ es el último número y la
razón entre números consecutivos es $k$.
Por ejemplo,
\[3+6+12+24=\frac{24 \cdot 2 - 3}{2-1} = 45.\]

Esta fórmula puede deducirse de la siguiente manera. Sea
\[ S = a + ak + ak^2 + \cdots + b .\]
Multiplicando ambos lados por $k$, obtenemos
\[ kS = ak + ak^2 + ak^3 + \cdots + bk,\]
y resolviendo la ecuación
\[ kS-S = bk-a\]
se obtiene la fórmula.

Un caso especial de la suma de una progresión geométrica es la fórmula
\[1+2+4+8+\ldots+2^{n-1}=2^n-1.\]

\index{suma armónica}

Una \key{suma armónica} es una suma de la forma
\[ \sum_{x=1}^n \frac{1}{x} = 1+\frac{1}{2}+\frac{1}{3}+\ldots+\frac{1}{n}.\]

Una cota superior para una suma armónica es $\log_2(n)+1$.
En efecto, podemos
modificar cada término $1/k$ de modo que $k$ se convierta
en la potencia de dos más cercana que no exceda a $k$.
Por ejemplo, cuando $n=6$, podemos estimar
la suma de la siguiente manera:
\[ 1+\frac{1}{2}+\frac{1}{3}+\frac{1}{4}+\frac{1}{5}+\frac{1}{6} \le
1+\frac{1}{2}+\frac{1}{2}+\frac{1}{4}+\frac{1}{4}+\frac{1}{4}.\]
Esta cota superior consta de $\log_2(n)+1$ partes
($1$, $2 \cdot 1/2$, $4 \cdot 1/4$, etc.),
y el valor de cada parte es a lo sumo 1.

\subsubsection{Teoría de conjuntos}

\index{teoría de conjuntos}
\index{conjunto}
\index{intersección}
\index{unión}
\index{diferencia}
\index{subconjunto}
\index{conjunto universal}
\index{complemento}

Un \key{conjunto} es una colección de elementos.
Por ejemplo, el conjunto
\[X=\{2,4,7\}\]
contiene los elementos 2, 4 y 7.
El símbolo $\emptyset$ denota un conjunto vacío,
y $|S|$ denota el tamaño de un conjunto $S$,
es decir, el número de elementos del conjunto.
Por ejemplo, en el conjunto anterior, $|X|=3$.

Si un conjunto $S$ contiene un elemento $x$,
escribimos $x \in S$,
y en caso contrario escribimos $x \notin S$.
Por ejemplo, en el conjunto anterior
\[4 \in X \hspace{10px}\textrm{y}\hspace{10px} 5 \notin X.\]

\begin{samepage}
Pueden construirse nuevos conjuntos usando operaciones de conjuntos:
\begin{itemize}
\item La \key{intersección} $A \cap B$ consta de los elementos
que están tanto en $A$ como en $B$.
Por ejemplo, si $A=\{1,2,5\}$ y $B=\{2,4\}$,
entonces $A \cap B = \{2\}$.
\item La \key{unión} $A \cup B$ consta de los elementos
que están en $A$ o en $B$ o en ambos.
Por ejemplo, si $A=\{3,7\}$ y $B=\{2,3,8\}$,
entonces $A \cup B = \{2,3,7,8\}$.
\item El \key{complemento} $\bar A$ consta de los elementos
que no están en $A$.
La interpretación de un complemento depende del
\key{conjunto universal}, que contiene todos los elementos posibles.
Por ejemplo, si $A=\{1,2,5,7\}$ y el conjunto universal es
$\{1,2,\ldots,10\}$, entonces $\bar A = \{3,4,6,8,9,10\}$.
\item La \key{diferencia} $A \setminus B = A \cap \bar B$
consta de los elementos que están en $A$ pero no en $B$.
Nota que $B$ puede contener elementos que no están en $A$.
Por ejemplo, si $A=\{2,3,7,8\}$ y $B=\{3,5,8\}$,
entonces $A \setminus B = \{2,7\}$.
\end{itemize}
\end{samepage}

Si cada elemento de $A$ también pertenece a $S$,
decimos que $A$ es un \key{subconjunto} de $S$,
denotado por $A \subset S$.
Un conjunto $S$ siempre tiene $2^{|S|}$ subconjuntos,
incluido el conjunto vacío.
Por ejemplo, los subconjuntos del conjunto $\{2,4,7\}$ son
\begin{center}
$\emptyset$,
$\{2\}$, $\{4\}$, $\{7\}$, $\{2,4\}$, $\{2,7\}$, $\{4,7\}$ y $\{2,4,7\}$.
\end{center}

Algunos conjuntos de uso frecuente son
$\mathbb{N}$ (números naturales),
$\mathbb{Z}$ (enteros),
$\mathbb{Q}$ (números racionales) y
$\mathbb{R}$ (números reales).
El conjunto $\mathbb{N}$
puede definirse de dos maneras, según
la situación:
o bien $\mathbb{N}=\{0,1,2,\ldots\}$
o bien $\mathbb{N}=\{1,2,3,...\}$.

También podemos construir un conjunto usando una regla de la forma
\[\{f(n) : n \in S\},\]
donde $f(n)$ es alguna función.
Este conjunto contiene todos los elementos de la forma $f(n)$,
donde $n$ es un elemento de $S$.
Por ejemplo, el conjunto
\[X=\{2n : n \in \mathbb{Z}\}\]
contiene todos los enteros pares.

\subsubsection{Lógica}

\index{lógica}
\index{negación}
\index{conjunción}
\index{disyunción}
\index{implicación}
\index{equivalencia}

El valor de una expresión lógica es
\key{verdadero} (1) o \key{falso} (0).
Los operadores lógicos más importantes son
$\lnot$ (\key{negación}),
$\land$ (\key{conjunción}),
$\lor$ (\key{disyunción}),
$\Rightarrow$ (\key{implicación}) y
$\Leftrightarrow$ (\key{equivalencia}).
La siguiente tabla muestra los significados de estos operadores:

\begin{center}
\begin{tabular}{rr|rrrrrrr}
$A$ & $B$ & $\lnot A$ & $\lnot B$ & $A \land B$ & $A \lor B$ & $A \Rightarrow B$ & $A \Leftrightarrow B$ \\
\hline
0 & 0 & 1 & 1 & 0 & 0 & 1 & 1 \\
0 & 1 & 1 & 0 & 0 & 1 & 1 & 0 \\
1 & 0 & 0 & 1 & 0 & 1 & 0 & 0 \\
1 & 1 & 0 & 0 & 1 & 1 & 1 & 1 \\
\end{tabular}
\end{center}

La expresión $\lnot A$ tiene el valor opuesto de $A$.
La expresión $A \land B$ es verdadera si tanto $A$ como $B$
son verdaderas,
y la expresión $A \lor B$ es verdadera si $A$ o $B$ o ambas
son verdaderas.
La expresión $A \Rightarrow B$ es verdadera
si siempre que $A$ es verdadera, también $B$ es verdadera.
La expresión $A \Leftrightarrow B$ es verdadera
si $A$ y $B$ son ambas verdaderas o ambas falsas.

\index{predicado}

Un \key{predicado} es una expresión que es verdadera o falsa
según sus parámetros.
Los predicados suelen denotarse con letras mayúsculas.
Por ejemplo, podemos definir un predicado $P(x)$
que es verdadero exactamente cuando $x$ es un número primo.
Usando esta definición, $P(7)$ es verdadero pero $P(8)$ es falso.

\index{cuantificador}

Un \key{cuantificador} conecta una expresión lógica
con los elementos de un conjunto.
Los cuantificadores más importantes son
$\forall$ (\key{para todo}) y $\exists$ (\key{existe}).
Por ejemplo,
\[\forall x (\exists y (y < x))\]
significa que para cada elemento $x$ del conjunto,
existe un elemento $y$ en el conjunto
tal que $y$ es menor que $x$.
Esto es verdadero en el conjunto de los enteros,
pero falso en el conjunto de los números naturales.

Usando la notación descrita arriba,
podemos expresar muchos tipos de proposiciones lógicas.
Por ejemplo,
\[\forall x ((x>1 \land \lnot P(x)) \Rightarrow (\exists a (\exists b (a > 1 \land b > 1 \land x = ab))))\]
significa que si un número $x$ es mayor que 1
y no es un número primo,
entonces existen números $a$ y $b$
que son mayores que $1$ y cuyo producto es $x$.
Esta proposición es verdadera en el conjunto de los enteros.

\subsubsection{Funciones}

La función $\lfloor x \rfloor$ redondea el número $x$
hacia abajo a un entero, y la función
$\lceil x \rceil$ redondea el número $x$
hacia arriba a un entero. Por ejemplo,
\[ \lfloor 3/2 \rfloor = 1 \hspace{10px} \textrm{y} \hspace{10px} \lceil 3/2 \rceil = 2.\]

Las funciones $\min(x_1,x_2,\ldots,x_n)$
y $\max(x_1,x_2,\ldots,x_n)$
dan el menor y el mayor de los valores
$x_1,x_2,\ldots,x_n$.
Por ejemplo,
\[ \min(1,2,3)=1 \hspace{10px} \textrm{y} \hspace{10px} \max(1,2,3)=3.\]

\index{factorial}

El \key{factorial} $n!$ puede definirse
\[\prod_{x=1}^n x = 1 \cdot 2 \cdot 3 \cdot \ldots \cdot n\]
o recursivamente
\[
\begin{array}{lcl}
0! & = & 1 \\
n! & = & n \cdot (n-1)! \\
\end{array}
\]

\index{número de Fibonacci}

Los \key{números de Fibonacci}
%\footnote{Fibonacci (c. 1175--1250) was an Italian mathematician.}
aparecen en muchas situaciones.
Pueden definirse recursivamente de la siguiente manera:
\[
\begin{array}{lcl}
f(0) & = & 0 \\
f(1) & = & 1 \\
f(n) & = & f(n-1)+f(n-2) \\
\end{array}
\]
Los primeros números de Fibonacci son
\[0, 1, 1, 2, 3, 5, 8, 13, 21, 34, 55, \ldots\]
También existe una fórmula cerrada
para calcular los números de Fibonacci, que a veces se llama
\index{fórmula de Binet} \key{fórmula de Binet}:
\[f(n)=\frac{(1 + \sqrt{5})^n - (1-\sqrt{5})^n}{2^n \sqrt{5}}.\]

\subsubsection{Logaritmos}

\index{logaritmo}

El \key{logaritmo} de un número $x$
se denota $\log_k(x)$, donde $k$ es la base
del logaritmo.
Según la definición,
$\log_k(x)=a$ exactamente cuando $k^a=x$.

Una propiedad útil de los logaritmos es
que $\log_k(x)$ es igual al número de veces
que tenemos que dividir $x$ entre $k$ antes de llegar
al número 1.
Por ejemplo, $\log_2(32)=5$
porque se necesitan 5 divisiones entre 2:

\[32 \rightarrow 16 \rightarrow 8 \rightarrow 4 \rightarrow 2 \rightarrow 1 \]

Los logaritmos se usan a menudo en el análisis de
algoritmos, porque muchos algoritmos eficientes
reducen algo a la mitad en cada paso.
Por eso, podemos estimar la eficiencia de tales algoritmos
usando logaritmos.

El logaritmo de un producto es
\[\log_k(ab) = \log_k(a)+\log_k(b),\]
y en consecuencia,
\[\log_k(x^n) = n \cdot \log_k(x).\]
Además, el logaritmo de un cociente es
\[\log_k\Big(\frac{a}{b}\Big) = \log_k(a)-\log_k(b).\]
Otra fórmula útil es
\[\log_u(x) = \frac{\log_k(x)}{\log_k(u)},\]
y usando esto, es posible calcular
logaritmos en cualquier base si hay una forma de
calcular logaritmos en alguna base fija.

\index{logaritmo natural}

El \key{logaritmo natural} $\ln(x)$ de un número $x$
es un logaritmo cuya base es $e \approx 2.71828$.
Otra propiedad de los logaritmos es que
el número de dígitos de un entero $x$ en base $b$ es
$\lfloor \log_b(x)+1 \rfloor$.
Por ejemplo, la representación de
$123$ en base $2$ es 1111011 y
$\lfloor \log_2(123)+1 \rfloor = 7$.

\section{Concursos y recursos}

\subsubsection{IOI}

La Olimpiada Internacional de Informática (IOI)
es un concurso de programación anual para
estudiantes de secundaria.
Cada país puede enviar un equipo de
cuatro estudiantes al concurso.
Suele haber unos 300 participantes
de 80 países.

La IOI consta de dos concursos de cinco horas de duración.
En ambos concursos, se pide a los participantes que
resuelvan tres tareas algorítmicas de dificultad variada.
Las tareas se dividen en subtareas,
cada una de las cuales tiene una puntuación asignada.
Aunque los concursantes se dividen en equipos,
compiten como individuos.

El temario de la IOI \cite{iois} regula los temas
que pueden aparecer en las tareas de la IOI.
Casi todos los temas del temario de la IOI
están cubiertos por este libro.

Los participantes de la IOI se seleccionan mediante
concursos nacionales.
Antes de la IOI, se organizan muchos concursos regionales,
como la Olimpiada Báltica de Informática (BOI),
la Olimpiada Centroeuropea de Informática (CEOI)
y la Olimpiada de Informática de Asia-Pacífico (APIO).

Algunos países organizan concursos de práctica en línea
para futuros participantes de la IOI,
como el Concurso Abierto Croata de Informática \cite{coci}
y la Olimpiada de Computación de EE.\ UU. \cite{usaco}.
Además, hay una gran colección de problemas de concursos polacos
disponible en línea \cite{main}.

\subsubsection{ICPC}

El Concurso Internacional de Programación Colegial (ICPC)
es un concurso de programación anual para estudiantes universitarios.
Cada equipo del concurso consta de tres estudiantes,
y a diferencia de la IOI, los estudiantes trabajan juntos;
solo hay una computadora disponible para cada equipo.

El ICPC consta de varias etapas, y finalmente los
mejores equipos son invitados a las Finales Mundiales.
Aunque hay decenas de miles de participantes
en el concurso, solo hay un pequeño número\footnote{El número exacto de plazas para la final
varía de año en año; en 2017, hubo 133 plazas para la final.} de plazas disponibles para la final,
así que incluso avanzar a las finales
es un gran logro en algunas regiones.

En cada concurso ICPC, los equipos tienen cinco horas de tiempo para
resolver unos diez problemas algorítmicos.
La solución de un problema se acepta solo si resuelve
todos los casos de prueba de forma eficiente.
Durante el concurso, los competidores pueden ver los resultados de otros equipos,
pero durante la última hora el marcador se congela y no
es posible ver los resultados de los últimos envíos.

Los temas que pueden aparecer en el ICPC no están tan bien
especificados como los de la IOI.
En cualquier caso, está claro que se necesita más conocimiento
en el ICPC, especialmente más habilidades matemáticas.

\subsubsection{Concursos en línea}

También hay muchos concursos en línea que están abiertos para todos.
En este momento, el sitio de concursos más activo es Codeforces,
que organiza concursos aproximadamente cada semana.
En Codeforces, los participantes se dividen en dos divisiones:
los principiantes compiten en la Div2 y los programadores más experimentados en la Div1.
Otros sitios de concursos son AtCoder, CS Academy, HackerRank y Topcoder.

Algunas empresas organizan concursos en línea con finales presenciales.
Ejemplos de tales concursos son el Facebook Hacker Cup,
el Google Code Jam y Yandex.Algorithm.
Por supuesto, las empresas también usan esos concursos para reclutar:
tener un buen desempeño en un concurso es una buena forma de demostrar las habilidades de uno.

\subsubsection{Libros}

Ya hay algunos libros (además de este) que
se centran en la programación competitiva y la resolución algorítmica de problemas:

\begin{itemize}
\item S. S. Skiena y M. A. Revilla:
\emph{Programming Challenges: The Programming Contest Training Manual} \cite{ski03}
\item S. Halim y F. Halim:
\emph{Competitive Programming 3: The New Lower Bound of Programming Contests} \cite{hal13}
\item K. Diks et al.: \emph{Looking for a Challenge? The Ultimate Problem Set from
the University of Warsaw Programming Competitions} \cite{dik12}
\end{itemize}

Los dos primeros libros están pensados para principiantes,
mientras que el último libro contiene material avanzado.

Por supuesto, los libros generales de algoritmos también son adecuados para
los programadores competitivos.
Algunos libros populares son:

\begin{itemize}
\item T. H. Cormen, C. E. Leiserson, R. L. Rivest y C. Stein:
\emph{Introduction to Algorithms} \cite{cor09}
\item J. Kleinberg y É. Tardos:
\emph{Algorithm Design} \cite{kle05}
\item S. S. Skiena:
\emph{The Algorithm Design Manual} \cite{ski08}
\end{itemize}
